\scanpage{38}\mspage{$\underline{\overline{\mathrm{XI}}}$}
3 - \uline{Preuve du théorĕme} 1

\margin{!}Soit $\Gamma$ un polygone et $E(\Gamma)$ son enveloppe convexe. Soit $s,s'\in S(\Gamma)$ un couple de sommets \deleted{\illegible{1}} tels que $s'$ est $(\mathbb{R}^2-\Gamma)$-visible de $s$ et $]s,s'[\subset\partial E(\Gamma)$ (cf. coroll. au lemme $\underline{\overline{\mathrm{VIII}}}$). Soient $\Gamma^1_{ss'}$ et $\Gamma^2_{ss'}$ les deux arcs de polygone \inserted{(de $\Gamma$)} réliant $s$ à $s'$ (cf. Remarque 1).
\margin{\tikz[baseline=0pt]{\draw (0,0) .. controls (-.04,.24) and (-.03,.52) .. (-.02,.65);}}

Soit $\Lambda_{ss'}$ l'arc de polygone $[s,s']$. On appelle $\Lambda^1$ (resp. $\Lambda^2$) le polygone obtenu en récollant $\Lambda_{ss'}$ avec $\Gamma^1_{ss'}$ (resp. $\Gamma^2_{ss'}$) selon $\{s,s'\}$.

% slanted
\margin{\shortstack{$a=s$\\$b=s'$}\,?}
\uline{Lemme} $\underline{\overline{\mathrm{XII}}}$. Si $\Lambda^1$ (resp. $\Lambda^2$) est un $J$-polygone alors $\Gamma^2_{ab}-\{a,b\}\subset\operatorname{int}\Lambda^1$ ou bien $\Gamma^2_{ab}-\{a,b\}\subset\operatorname{ext}\Lambda^{\deleted{\unsure{2}}1}$ (resp. $\Gamma^1_{ab}-\{a,b\}\subset\operatorname{int}\Lambda^2$ $\Gamma^1_{ab}-\{a,b\}\subset\operatorname{ext}\Lambda^2$)
\margin{1}

\uline{Demonstration} : Il est clair que $\Gamma^2_{ab}-\{a,b\}\subset\deleted{\unsure{E|\partial}}\mathbb{R}^2-\Lambda^1$ car $]s,s'[\cap\Gamma=\unsure{\phi|\varnothing}$, mais en plus $\Gamma^2_{ab}-\{a,b\}$ est connexe donc la conclusion. On réfait l'argument pour $\Lambda^2$, c.q.f.d.
\margin{\tikz[baseline=0pt]{\draw (0,0) .. controls (-.03,.20) and (-.02,.50) .. (-.03,.65);}}
\margin{\tikz[baseline=0pt]{\draw (0,0) -- (0,.65);}}

On suppose maintenant \inserted{que} $\Lambda^1$ et $\Lambda^2$ \deleted{\unsure{que}} sont des $J$-polygones.

\uline{Lemme} $\underline{\overline{\mathrm{XIII}}}$ : Si $x\in\partial E(\Gamma)$ et $D_x$ est un disque de centre $x$ \unsure{vérifiont|vérifiant} b) de la definition 2 par rapport à $\partial E(\Gamma)$ on a alors :
\[
\operatorname{ext}\partial E(\Gamma)\cap\text{\tikz[baseline=(dx.base)]{\node[inner sep=2pt] (dx) {$\deleted{D_x}$};\draw ([xshift=-2pt,yshift=2pt]dx.north west) -- ([xshift=2pt,yshift=-4pt]dx.south east);\draw ([xshift=2pt,yshift=4pt]dx.north east) -- ([xshift=-1pt,yshift=-4pt]dx.south west);}}\subset\operatorname{ext}\Lambda^i,\qquad i=1,2.
\]
\margin{\resizebox{\marginparwidth}{!}{\tikz[baseline=(ex.base)]{\node[inner sep=3pt] (ex) {$\begin{gathered}
\operatorname{Ext}(\partial E(\Gamma))=\mathbb{R}^{\unsure{2|7}}-E(\Gamma)\\
\subset\operatorname{Ext}(\Lambda^i)
\end{gathered}$};
\draw (ex.north east) to[bend right=3] (ex.north west) -- (ex.south west) to[bend left=3] (ex.south);}}}

\uline{Démonstration} : Comme $E(\Lambda^i)\subset E(\Gamma)$, $i=1,2$, on a $\operatorname{ext}\deleted{\partial E(}\Lambda^i\deleted{)}\supset\operatorname{ext}\partial E(\Lambda^i)\supset\operatorname{ext}\partial E(\Lambda)$ ce qui entraîne l'ennoncé.
% archive sheet number: 37
